À partir de la solution précédente $\mathrm{(S)}$, on prépare une solution
aqueuse $\mathrm{(S_{A})}$ de concentration molaire $C_{A}$. Le taux
d'avancement final de la réaction entre l'acide méthanoïque et l'eau dans ce cas
vaut $\tau=0,21$.

\begin{donnees}
    \begin{itemize}
        \item $\mathrm{pK_{A}}(\ce{HCO2H_{(aq)}}/\ce{HCO^-_{2(aq)}})=3,8$
    \end{itemize}
\end{donnees}
\begin{enumerate}
    \item Écrire l'équation chimique de la réaction de l'acide méthanoïque avec
          l'eau.
    \item Montrer que $C_{A}=\frac{1-\tau}{\tau^{2}}\cdot10^{-\mathrm{pK_{A}}}$.
          Calculer la valeur de $C_{A}$.
    \item Déterminer la valeur du $\mathrm{pH}$ de la solution
          $\mathrm{(S_{A})}$.
    \item En utilisant le dispositif schématisé sur la
          Fig.~\ref{fig:2bac_svt_national_2023_rattrapage_chim_1}, on dose le
          volume $V_{A}=\qty{30}{\mL}$ de la solution $\mathrm{(S_{A})}$ par une
          solution aqueuse $\mathrm{(S_{B})}$ d'hydroxyde de sodium
          $\ce{Na^+_{(aq)} + HO^-_{(aq)}}$de concentration molaire
          $C_{B}=\qty{1e-2}{\mol\per\L}$. La courbe de la
          Fig.~\ref{fig:2bac_svt_national_2023_rattrapage_chim_2}
          représente la variation du $\mathrm{pH}$ du mélange en fonction du
          volume $V_{B}$ de la solution $\mathrm{(S_{B})}$ versé au cours du
          dosage.

          \begin{minipage}{0.49\linewidth}
              \begin{figure}[H]
                  \centering
                  \includegraphics[width=0.8\textwidth]{2026-03-06_115446-}
                  \caption{}
                  \label{fig:2bac_svt_national_2023_rattrapage_chim_1}
              \end{figure}
          \end{minipage}
          \hfill
          \begin{minipage}{0.49\linewidth}
              \begin{figure}[H]
                  \centering
                  \includegraphics[width=0.9\textwidth]{2026-03-06_115432-}
                  \caption{}
                  \label{fig:2bac_svt_national_2023_rattrapage_chim_2}
              \end{figure}
          \end{minipage}
          
          \begin{enumerate}
              \item Nommer les éléments du dispositif numérotés 1, 2 et 3.
              \item Déterminer les coordonnées du point d'équivalence.
              \item Retrouver la valeur de $C_{A}$.
              \item Préciser, parmi les indicateurs colorés ci-dessous,
                    l'indicateur convenable pour repérer cette équivalence.
          
                    \begin{minipage}{\linewidth}
                        \begin{table}[H]
                            \centering
                            \begin{tabularx}{\textwidth}{|p{4cm}|C|C|C|}
                                \hline
                                \textbf{Indicateur coloré} &
                                \textbf{Teinte acide} &
                                \textbf{Zone de virage} &
                                \textbf{Teinte basique} \\
                                \hline
                                Rouge de méthyle & Rouge & $4,4-6,2$ & Jaune \\
                                \hline
                                Rouge de crésol & Jaune & $7,2-8,8$ & Rouge \\
                                \hline
                                Alizarine & Rouge & $11,0-12,4$ & Violette \\
                                \hline
                            \end{tabularx}
                        \end{table}
                    \end{minipage}
          \end{enumerate}
\end{enumerate}

